\documentclass[11pt,a4paper]{article}
\usepackage[T1]{fontenc}
\usepackage{lmodern,microtype}
\usepackage[margin=27mm]{geometry}
\usepackage{amsmath,amssymb,amsthm,mathtools}
\usepackage[authoryear,round]{natbib}
\usepackage{enumitem,booktabs,longtable}
\usepackage{xcolor}
\usepackage[colorlinks=true,linkcolor=blue!45!black,citecolor=blue!45!black,urlcolor=blue!45!black]{hyperref}
\usepackage{url}
\emergencystretch=1.5em
\setlength{\parskip}{3pt}
\setlength{\parindent}{1.2em}
\newcommand{\Q}{\mathbb Q}
\newcommand{\Z}{\mathbb Z}
\newcommand{\F}{\mathbb F}
\newcommand{\A}{\mathbb A}
\newcommand{\R}{\mathbb R}
\newcommand{\C}{\mathbb C}
\newcommand{\G}{\mathbb G}
\newcommand{\Zhat}{\widehat{\Z}}
\newcommand{\Qbar}{\overline{\Q}}
\newcommand{\kbar}{\overline{k}}
\newcommand{\eps}{\varepsilon}
\newcommand{\et}{\'et}
\newcommand{\ab}{\mathrm{ab}}
\newcommand{\fcov}{\mathrm{f\text{-}cov}}
\newcommand{\fab}{\mathrm{f\text{-}ab}}
\newcommand{\Sel}{\mathrm{Sel}}
\DeclareMathOperator{\Gal}{Gal}
\DeclareMathOperator{\GL}{GL}
\DeclareMathOperator{\SL}{SL}
\DeclareMathOperator{\Spec}{Spec}
\DeclareMathOperator{\Sec}{Sec}
\DeclareMathOperator{\Jac}{Jac}
\DeclareMathOperator{\Hom}{Hom}
\DeclareMathOperator{\im}{im}
\DeclareMathOperator{\ord}{ord}
\newcommand{\workbenchurl}{https://github.com/KokunoYumeto/k-theory-and-galois-groups-workbench}
\newcommand{\recordurl}{https://github.com/KokunoYumeto/mathematics-workbenches}
\theoremstyle{plain}
\newtheorem{theorem}{Theorem}[section]
\newtheorem{proposition}[theorem]{Proposition}
\newtheorem{lemma}[theorem]{Lemma}
\newtheorem{corollary}[theorem]{Corollary}
\theoremstyle{definition}
\newtheorem{definition}[theorem]{Definition}
\newtheorem{remark}[theorem]{Remark}
\newtheorem{hypothesis}[theorem]{Hypothesis}
\title{Fermat's theorem for Selmer sections\\[4pt]
\large Grothendieck's question, after Steklov}
\author{}
\date{25 November 2026}
\begin{document}
\maketitle
\begin{abstract}
Grothendieck proposed Fermat's theorem as a test of his section conjecture, and
Steklov adapted the modular method to locally geometric sections of prime-exponent
Fermat curves. We prove that, for every odd exponent at least three, every section
which is locally a point section or a rational cusp section is globally cuspidal.
The section set is exactly three disjoint cusp packets. The proof combines the
refined theorem of Serre, Stoll's finite descent theorems, and Stix's description
of cusp packets. It keeps the constant term in the cuspidal Tate parameter and
uses a modular curve of level four times the exponent; these changes also cover
exponents five and seven. Exponent three follows from the modular elliptic curve
of conductor twenty-seven, and finite coverings give every odd exponent.
We separate the statement about all sections, whose additional input is the
proper local section theorem asserted in OpenAI's unrefereed preprint. We also
treat positive-genus modular curves with a marked torsion point and the real
obstruction for even Fermat exponents. The mathematics and text were produced
by GPT-6 Astra Pro, with separate same-model proof checks and earlier
Claude Opus 5.5 checks recorded below; no human has checked this paper.
\end{abstract}

\section{Grothendieck's question and the results}

This paper determines the locally geometric part of the section set of an
affine Fermat curve. Grothendieck's discussion in
\emph{La Longue Marche \`a travers la th\'eorie de Galois}, \S51.16,
asks for Fermat's theorem in the language of sections
\citep[\S51.16, fols.~733--736v, especially fol.~734]{grothendieck}. The point is to recognize the evident boundary
sections directly in the arithmetic fundamental group. Fermat's Last Theorem
removes rational points from the affine curve; the section conjecture predicts
that the boundary accounts for everything left.

For an integer $n\geq3$, let
\[
 F_n:\ X^n+Y^n+Z^n=0\ \subset\mathbf P^2_{\Q},\qquad
 U_n=F_n\setminus D_n,\qquad D_n=F_n\cap\{XYZ=0\}.
\]
The curve $F_n$ is smooth of genus $(n-1)(n-2)/2$, and $D_n$ has $3n$
geometric points. In particular $U_n$ is hyperbolic. Write
$G_k=\Gal(\kbar/k)$ for the absolute Galois group of a field $k$ of
characteristic zero. A geometric base point gives an exact sequence
\begin{equation}\label{eq:fundamental}
 1\longrightarrow\pi_1(U_{n,\Qbar})\longrightarrow\pi_1(U_n)
 \longrightarrow G_{\Q}\longrightarrow1.
\end{equation}
Every section in this paper is continuous, and sections are identified up to
conjugation by the geometric fundamental group.

A rational point gives a section by functoriality. A rational boundary point
$y$ gives a decomposition group $D_{\widetilde y}$, with exact sequence
\begin{equation}\label{eq:cusp}
 1\longrightarrow I_{\widetilde y}\simeq\Zhat(1)
 \longrightarrow D_{\widetilde y}\longrightarrow G_k\longrightarrow1.
\end{equation}
A section factoring through this group is \emph{cuspidal at $y$}. Its packet
is the set of all such section classes. A nonzero rational tangent vector
splits \eqref{eq:cusp}; arbitrary cuspidal sections need not come from an
individual tangent vector.

\begin{definition}\label{def:selmer}
A section of a smooth affine hyperbolic curve $U/k$, with smooth proper
compactification $X$, is \emph{Selmer} if at every finite place $v$ its
restriction is the section of a point of $U(k_v)$ or a section in a
decomposition group of a point of $(X\setminus U)(k_v)$. At a real place we
require a real point section or a real cusp section. There is no additional
condition at complex places. This definition uses the full profinite groups.
\end{definition}

\begin{theorem}[Fermat's theorem for Selmer sections]\label{thm:fermat}
Let $n\geq3$ be odd. Every Selmer section of \eqref{eq:fundamental} is
cuspidal. More precisely, if
\[
 y_1=(0:-1:1),\qquad y_2=(-1:0:1),\qquad y_3=(1:-1:0),
\]
then
\begin{equation}\label{eq:packets}
 \Sec(U_n/\Q)^{\Sel}=\mathcal P_{y_1}\ \coprod\
 \mathcal P_{y_2}\ \coprod\ \mathcal P_{y_3}.
\end{equation}
Each packet is naturally a torsor under
\[
 H^1(\Q,\Zhat(1))\simeq\varprojlim_m\Q^\times/(\Q^\times)^m.
\]
The inputs are the refined form of Serre's modularity theorem,
Stoll's Theorem~8.2 and Corollary~8.8 with his 2017 corrections,
Mazur's torsion theorem, and Stix's cusp-packet theorem. No local section
conjecture is an input to this assertion.
\end{theorem}

Thus the local point-or-cusp condition is enough to prove the prediction
Grothendieck attached to Fermat's theorem. The statement concerns all
cuspidal sections, including the entire Kummer completion of a tangent
parameter. It does not identify each packet with just $\Q^\times$.

The local assertion needed to remove the Selmer hypothesis is stated precisely
because it is logically separate from the global proof.

\begin{hypothesis}[Proper local section input]\label{hyp:local}
For every prime $p$, every finite extension $k/\Q_p$, and every smooth proper
geometrically connected curve $X/k$ of genus at least two, the point map
\[
 X(k)\longrightarrow\Sec(X/k)
\]
is bijective for the full profinite arithmetic fundamental group.
\end{hypothesis}

This is the statement of \citet[Theorem~1.1, p.~2]{openai}, in its
6 October 2026 version. That preprint is unrefereed; we use its statement as
an explicitly identified input and do not certify its proof here.

\begin{corollary}[All sections, with its local input]\label{cor:all}
Assume Hypothesis~\ref{hyp:local}. For every odd $n\geq3$, every section of
\eqref{eq:fundamental} is Selmer and therefore cuspidal. Its entire section
set is the disjoint union \eqref{eq:packets}. The only further local input
is the established real section theorem.
\end{corollary}

\begin{theorem}[Two consequences of finite descent and the real place]
\label{thm:other}
If the smooth modular curve $X_1(N)/\Q$ has positive genus, every Selmer
section of $Y_1(N)/\Q$ is cuspidal at a rational cusp. Under
Hypothesis~\ref{hyp:local}, the same assertion holds for every full section.
For every even integer $n\geq2$, $U_n/\Q$ has no Selmer section and, in
fact, no full section. The latter assertion needs only the real section
theorem, not Hypothesis~\ref{hyp:local}.
\end{theorem}

The modular-curve conclusion includes every level of genus at least two,
as well as the positive-genus levels excluded from the list of possible
rational torsion orders. It is a consequence of Stoll's theorem and
Mazur's theorem, with the passage from finite support to cusp sections
made explicit below.

\subsection{Earlier results and the division of credit}

Stoll proves that finite abelian descent cuts out the rational points on
every positive-genus modular curve of the indicated types
\citep[Corollary~8.8, pp.~385--386]{stoll}. His density-one finite-support
theorem \citep[Theorem~8.2, pp.~383--384]{stoll} is what makes that result
useful for section neighborhoods. Stix proves the disjointness and internal
description of cusp packets \citep[Proposition~8 and Theorem~14,
pp.~5--6]{stix}. These are independent arithmetic and anabelian inputs.

Steklov associates a residual Legendre representation to a Fermat section,
uses modularity in the irreducible case, and finite descent on a modular
cover in the reducible case \citep[Theorem~A; \S5]{steklov}. This is the
proof strategy followed here. His printed statement has prime exponent
$\ell\geq7$. Two steps require correction: the asserted divisibility of
the cuspidal Tate class in Proposition~4.6, used in Proposition~5.10,
omits its leading constant; and the final appeal to
$Y_1(7)(\Q)=\varnothing$ in Theorem~5.14 is false. Section~\ref{sec:corrections}
gives the precise corrections. The local assertion required in
Proposition~5.10 remains true, so the first defect does not destroy the
conclusion for $\ell\geq11$ once that assertion is reproved. Exponent seven
is unproved by the printed argument. The map to $X_0(4\ell)$ below and
the full rational two-torsion give both five and seven.

\citet[Theorems~A, C, and E]{bkl} relate locally geometric sections to
finite descent and to Chabauty--Kim loci. Their affine integral conditions
are different from allowing all cusp packets at every place. Their
Theorem~E gives a finite algebraic support for Chabauty--Kim loci when a
Jacobian has a rank-zero factor of dimension at least two with finite
Tate--Shafarevich group; it does not by itself identify that support with
rational cusps. The companion paper discusses this distinction for the
Fermat quintic and for quotients of the geometric fundamental group.

\citet[Theorems~1.1--1.2, pp.~1--2, and Theorem~2.2,
pp.~3--4]{porowski} proves that a noncuspidal Selmer section is
determined by its reductions on a density-one set of primes, and that
a section which is locally cuspidal on such a set is globally cuspidal.
These are results for the full fundamental group. They clarify why
controlling the reductions of a section is a central step; the modular
argument here produces the needed global support for Fermat curves.

\subsection{Proof overview}

The main difficulty is that a section is not a rational point. One cannot
start by attaching an elliptic curve over $\Q$ to it. Instead the Legendre
family gives a two-dimensional representation of the fundamental group,
and composing with the section gives a Galois representation. Local
geometry makes its twists satisfy the same small-conductor and
finite-flat conditions that occur in the modular proof of Fermat's theorem.
The constant in a Tate parameter must be kept at a cusp: ramification
makes the variable part disappear modulo the exponent, but it leaves the
constant $2^{-8}$.

For a prime exponent $\ell\geq5$, a global quadratic twist of this
representation has conductor two and Serre weight two. If it were
irreducible, refined modularity would produce a weight-two cusp form of
level two, which does not exist. An invariant line therefore gives a
finite \et ale neighborhood of the section. The Legendre family's two
independent rational two-torsion points turn its cyclic $\ell$-subgroup
into a cyclic $4\ell$-subgroup on a two-isogenous curve. This produces a
nonconstant map to $X_0(4\ell)$, whose genus is positive even for $\ell=5$.

Stoll then forces the local points of the neighborhood into one finite
fiber. Applying the same theorem throughout the section's neighborhood
tower gives a compatible rational point in its compactified tower. Its
stabilizer is a point or cusp decomposition group. Rational points are
excluded using the same local representation, together with Mazur's
torsion list. The cubic is handled by $X_0(27)$. A coordinate power map
reduces every composite odd exponent to a prime exponent. Finally Stix's
theorem identifies the packets without identifying any two of them.

The Selmer condition is used twice: to control every local residual
representation, and to place the compactification points in finite
descent sets. Oddness is used to obtain the three rational cusps and in
the coordinate-power reduction. The bound $\ell\geq5$ makes the
two-adic $j$-invariant negative and $X_0(4\ell)$ of positive genus.
Neither step is asserted for the cubic.

\section{Finite support and section neighborhoods}\label{sec:finite}

This section spells out the passage from finite adelic support to an
arithmetic decomposition group. It is the mechanism of
\citet[Lemma~3.16 and Proposition~3.17]{steklov}, based on Stoll's
theorem. We give the argument in the form needed here, including its
behavior at a ramified boundary point and at real places.

For a smooth proper curve $X/k$ over a number field, set
\[
 X(\A_k)_\bullet=
 \prod_{v\text{ finite}}X(k_v)\times
 \prod_{v\text{ real}}\pi_0(X(k_v)).
\]
Complex factors are omitted, as required by
\citet[\S1]{stollerrata}. The subsets
$X(\A_k)_\bullet^{\fcov}$ and
$X(\A_k)_\bullet^{\fab}$ consist of points surviving all torsors under
finite \et ale $k$-group schemes, respectively all such abelian torsors.
Survival means that some global twist of the torsor has local points
above the specified adelic point at every place. In particular
\begin{equation}\label{eq:descentinclusion}
 X(k)\subset X(\A_k)_\bullet^{\fcov}
 \subset X(\A_k)_\bullet^{\fab}.
\end{equation}

\begin{lemma}[Lifting local geometry]\label{lem:lifting}
Let $U/k$ be a smooth geometrically connected affine curve, and let $s$
be a Selmer section. A connected finite \et ale cover $V\to U$ to which
$s$ lifts is geometrically connected, and the lifted section is Selmer.
If $X_V$ is its smooth proper compactification, the local point-or-cusp
centers give a point of $X_V(\A_k)_\bullet^{\fcov}$.
\end{lemma}
\begin{proof}
Let $\Pi=\pi_1(U)$, $\Delta=\pi_1(U_{\kbar})$, and let the open subgroup
$H\subset\Pi$ define $V$. Choose the lift so that $s(G_k)\subset H$.
The image of $H$ in $G_k$ is therefore all of $G_k$, which is equivalent
to geometric connectedness of $V$.

Fix a place $v$. If $s_v$ is the section of an interior point, choose a
prolongation whose decomposition group is $s_v(G_{k_v})$. That group
lies in $H$, so its image on $V$ is a $k_v$-point and induces the lifted
section. If $s_v$ is cuspidal, choose a prolongation $\widetilde c$
whose decomposition group contains $s_v(G_{k_v})$. Its intersection
with $H$ still maps onto $G_{k_v}$, because it contains the section.
Consequently the boundary point below $\widetilde c$ on $X_V$ is
$k_v$-rational, and the lifted section lies in its decomposition group.
The same argument applies over $\R$.

Compose the lifted section with $\pi_1(V)\to\pi_1(X_V)$. At a boundary
point this kills inertia, so locally it is the point section of the
center just obtained. For a torsor $T\to X_V$ under a finite \et ale
$k$-group scheme $G$, pull its cohomology class back along this global
section. The result is a single class $\alpha\in H^1(k,G)$. At every
place its restriction is the class of the fiber over the local center.
Twisting by this class, with the inverse convention if needed, makes
all these fibers trivial. This proves survival of every finite torsor.
Properness supplies the adelic condition at all finite places.
\end{proof}

\begin{lemma}[Finite support]\label{lem:finite-support}
Let $U/k$ be a smooth geometrically connected affine curve whose smooth
proper compactification $X$ has genus at least one. Let $s$ be Selmer.
Suppose there is a finite reduced subscheme $Z\subset X$ such that its
local centers belong to $Z(k_v)$ at a density-one set of finite places.
Then $s$ is the section of a point of $U(k)$, or is cuspidal at a point
of $(X\setminus U)(k)$.
\end{lemma}
\begin{proof}
The geometric group $\Delta$ is topologically finitely generated.
For every characteristic open subgroup $N\subset\Delta$, put
$H_N=N\,s(G_k)$. This is an open subgroup of $\Pi$, and the corresponding
cover $U_N\to U$ is geometrically connected. These covers are cofinal
among the finite section neighborhoods: given an open subgroup
containing $s(G_k)$, take a characteristic open subgroup inside its
intersection with $\Delta$. Their compactifications $X_N$ have genus
at least one by Riemann--Hurwitz.

By Lemma~\ref{lem:lifting}, $X_N$ has an adelic point surviving finite
descent. At the specified density-one set, its centers lie in the finite
inverse image $Z_N$ of $Z$. Stoll's Theorem~8.2 therefore gives a
$k$-point of $Z_N$. This application uses only finite places. In
particular it makes no claim that a chosen real connected component
has a point in a prescribed scheme-theoretic fiber; that is the issue
corrected in \citet[\S3]{stollerrata}.

The finite sets $Z_N(k)$ are nonempty and form an inverse system under
the transition morphisms. Its inverse limit is nonempty. Indeed, for
any finite collection of levels, a common refinement supplies a point
compatible with all transition conditions involving those levels;
compactness of the product of the finite sets then gives a point
$(z_N)_N$ in the inverse limit. Its image $z$ on $X$ is rational.

For completeness, interpret this limit in the geometric cover tower.
The section $s$ identifies $\Pi$ with $\Delta\rtimes G_k$. The geometric
cover with group $N$ acquires its $k$-structure from this action; its
compactified fiber above $z$ consists of the $N$-orbits of geometric
prolongations of $z$. A rational $z_N$ is an $s(G_k)$-fixed such orbit.
A compatible system of these orbits is a prolongation $\widetilde z$
in the full compactified geometric tower, fixed by $s(G_k)$. One may
also see the last assertion by taking the inverse limit of the finite
orbit sets: the action commutes with every transition map, and
characteristic open subgroups form a basis at the identity.

The stabilizer of $\widetilde z$ in $\Pi$ is its decomposition group
$D_{\widetilde z}$. Thus
$s(G_k)\subset D_{\widetilde z}$. If $z\in U(k)$, this decomposition
group maps isomorphically to $G_k$ and the section is the point section.
If $z\in X\setminus U$, it is a cusp section. At a cusp the conclusion
is containment of the section in the decomposition group; the
nontrivial cusp inertia is not contained in the section.
\end{proof}

\begin{corollary}[A map to a curve controlled by descent]\label{cor:descent-map}
In the setting of Lemma~\ref{lem:finite-support}, suppose there is a
nonconstant morphism $f:X\to W$ to a smooth proper curve such that
\[
 W(\A_k)_\bullet^{\fab}=W(k).
\]
Then every Selmer section of $U$ is a point section or a cusp section.
\end{corollary}
\begin{proof}
The centers in $X(\A_k)_\bullet^{\fcov}$ from
Lemma~\ref{lem:lifting} map into $W(\A_k)_\bullet^{\fab}$: pull back
every torsor on $W$. They therefore map to a single point $w\in W(k)$.
At each finite place their centers lie in the finite fiber $f^{-1}(w)$.
Apply Lemma~\ref{lem:finite-support} to its reduction.
\end{proof}

\section{The Tate parameter at the boundary}\label{sec:cusp}

The local calculation in this section retains the constant that ramification
alone does not control. Let
$T=\mathbf P^1\setminus\{0,1,\infty\}$ and let
\[
 E_\lambda:\quad y^2=x(x-1)(x-\lambda)
\]
be its Legendre elliptic scheme. Its invariants are
\begin{equation}\label{eq:legendre}
 \begin{split}
 \Delta(E_\lambda)&=16\lambda^2(1-\lambda)^2,\\
 c_4(E_\lambda)&=16(\lambda^2-\lambda+1),\\
 j(\lambda)&=256\frac{(\lambda^2-\lambda+1)^3}
                         {\lambda^2(1-\lambda)^2}.
 \end{split}
\end{equation}
The morphism $U_n\to T$ used throughout is
$\lambda=-(X/Z)^n$; the equation of $F_n$ gives
$1-\lambda=-(Y/Z)^n$.

\begin{lemma}[The constant in a formal Kummer class]\label{lem:kummer}
Let $k$ have characteristic zero, let $t$ be a parameter at a rational
boundary point, and suppose $q=a t^e h(t)$ with $a\in k^\times$,
$e\in\Z$, and $h(t)\in1+t k[[t]]$. For any section $s$ of the
punctured local fundamental group,
\begin{equation}\label{eq:kummer}
 \kappa(q)(s)=\kappa(a)+e\,\kappa(t)(s)
       \quad\text{in }H^1(k,\Zhat(1)).
\end{equation}
Here $\kappa$ denotes the multiplicative Kummer map, written additively.
\end{lemma}
\begin{proof}
For every positive integer $m$, the binomial expansion, equivalently
Hensel's lemma, gives an $m$th root of $h(t)$ in $1+t k[[t]]$.
Consequently $[q]=[a]+e[t]$ in the $m$th Kummer quotient.
Pullback along $s$ and passage to the compatible inverse limit prove
\eqref{eq:kummer}. The argument is equally valid in the punctured
henselian local scheme, whose finite tame covers in characteristic
zero have the same parameter and residue-field description.
\end{proof}

\begin{proposition}[Cuspidal Legendre class]\label{prop:cusp-class}
At every geometric cusp of $F_n$, over any characteristic-zero field
over which that cusp is rational, a parameter $t$ can be chosen so that
\begin{equation}\label{eq:jinverse}
 j^{-1}=2^{-8}t^{2n}
       \frac{(1+t^n)^2}{(1+t^n+t^{2n})^3}.
\end{equation}
The Tate parameter satisfies $q=j^{-1}+O(j^{-2})$. Hence every cusp
section satisfies
\[
 \kappa(q)(s)=\kappa(2^{-8})+2n\,\kappa(t)(s).
\]
In particular, for $n=\ell$ prime,
\begin{equation}\label{eq:cuspresidual}
 [q(s)]_\ell=[2^{-8}]_\ell\in H^1(k,\mu_\ell).
\end{equation}
The residual Legendre representation is a quadratic twist of the Tate
extension with this class; the extension need not split.
\end{proposition}
\begin{proof}
At a cusp with $X=0$, take $t=X/Z$, so $\lambda=-t^n$.
At a cusp with $Y=0$, take $t=Y/Z$, so $\lambda=1+t^n$.
At a cusp with $Z=0$, take $t=Z/X$, so $\lambda=-t^{-n}$.
Substitution in \eqref{eq:legendre} gives \eqref{eq:jinverse} in all
three cases. This does not require the other coordinate to be rational
over the original field: the formula holds over the residue field of
the chosen cusp.

The inverse of the integral Tate $j$-series has leading term $j^{-1}$.
Thus $q=2^{-8}t^{2n}h(t)$ with $h(0)=1$. Apply
Lemma~\ref{lem:kummer}. The standard Tate uniformization identifies
the $\ell$-torsion with an extension of $\Z/\ell$ by $\mu_\ell$
classified by $[q]$, after the quadratic character that makes the
reduction split multiplicative. Pulling back along the cusp section
gives the stated extension and its possible quadratic twist.
\end{proof}

\begin{remark}[The cusp above infinity]\label{rem:infinity}
The quadratic twist at infinity is not independent of the section.
For the Legendre equation,
$c_6=32(\lambda+1)(\lambda-2)(2\lambda-1)$, and with
$\lambda=-t^{-\ell}$ one has
\[
 -c_4/c_6=\tfrac14t^\ell(1+O(t^\ell)).
\]
Its square class is ramified in the parameter for odd $\ell$.
The local monodromy is $-T^{2\ell}$ rather than $T^{2\ell}$, where
$T$ is a unipotent transvection. Pulling back the quadratic character
along a cusp section can therefore depend on its tangent Kummer class.
The global patching below permits arbitrary local quadratic inertia
characters and includes precisely this possibility.
\end{remark}

The smallest relevant calculation already disproves divisibility by
ramification. With unit tangent at $X=0$ and $n=5$,
\[
 \kappa(q)(s)=-8\kappa(2).
\]
Its two-adic valuation is $-8$, whose image modulo five is $2$, not
zero. The order-ten ramification of the $j$-map does not remove that
constant.

\section{Local conditions and refined modularity}\label{sec:modularity}

This section proves that the residual representation of a Selmer section
has an invariant line. Fix a prime $\ell\geq5$, and let $s$ be a Selmer
section of $U_\ell$. Pull back the $\ell$-torsion of the Legendre scheme
and compose its monodromy with $s$:
\[
 \overline\rho_s:G_\Q\longrightarrow\GL_2(\F_\ell).
\]
The Weil pairing gives $\det\overline\rho_s=\overline\chi_\ell$, where
$\overline\chi_\ell$ is the mod-$\ell$ cyclotomic character.

\begin{lemma}[Local ordinary points]\label{lem:local-points}
Let $(a,b)\in U_\ell(\Q_p)$, so $a^\ell+b^\ell+1=0$, and set
$\lambda=-a^\ell$. At an odd prime $p$, the Legendre curve has good
reduction or is, after a quadratic twist, a Tate curve with
$\ell\mid v_p(q)$. At $p=2$ it is potentially multiplicative and
\[
 v_2(q)\equiv-8\pmod\ell.
\]
\end{lemma}
\begin{proof}
If $v_p(a)>0$, then $\lambda^2-\lambda+1$ and $1-\lambda$ are units,
so \eqref{eq:legendre} gives
\[
 v_p(j)=8v_p(2)-2\ell v_p(a).
\]
If $v_p(a)<0$, then
$v_p(1-\lambda)=v_p(\lambda)$ and
$v_p(\lambda^2-\lambda+1)=2v_p(\lambda)$; hence in both cases
\begin{equation}\label{eq:jvaluation}
 v_p(j)=8v_p(2)-2\ell|v_p(a)|.
\end{equation}
If $a$ is a unit but $\lambda\equiv1\pmod p$, use
$1-\lambda=-b^\ell$ instead. The remaining case at an odd prime has
$\lambda$ and $1-\lambda$ both units, and the discriminant in
\eqref{eq:legendre} is a unit, proving good reduction.

At two the remaining case cannot occur: if $a,b$ were both two-adic
units, $a^\ell+b^\ell+1$ would be odd. Thus
$v_2(j)=8-2\ell m$ for an integer $m\geq1$, which is negative because
$\ell\geq5$. At an odd prime every non-good case in
\eqref{eq:jvaluation} also has negative $j$-valuation. Tate
uniformization after a quadratic twist gives
$v_p(q)=-v_p(j)$, proving all the divisibility and congruence claims.
\end{proof}

\begin{lemma}[Local residual conditions]\label{lem:local-residual}
At every finite prime $p$ there is a quadratic character $\eps_p$ such
that $\overline\rho_s\otimes\eps_p$ has the following properties:
it is unramified if $p\neq2,\ell$; it is finite flat at $p=\ell$;
and at $p=2$ it is tamely ramified with conductor exponent one.
\end{lemma}
\begin{proof}
At an ordinary point choose the twist in Lemma~\ref{lem:local-points};
at a cusp use Proposition~\ref{prop:cusp-class}. In the good reduction
case the torsion of the smooth elliptic scheme has the asserted
unramified or finite-flat property, and no twist is needed.

In the Tate case the representation is
\[
 \begin{pmatrix}\overline\chi_\ell&c_q\\0&1\end{pmatrix},
 \qquad[c_q]=[q]\in H^1(\Q_p,\mu_\ell).
\]
For odd $p\neq\ell$ the valuation of the class is divisible by
$\ell$, both for points and for cusps, so it has a unit representative.
Taking $\ell$th roots is an automorphism on principal units. A root of
the residue-field unit lies in a finite residue-field extension and
lifts in the corresponding unramified extension. The Kummer class
and the cyclotomic character are therefore unramified.

At $p=\ell$ the same valuation condition gives a unit representative
$u\in\Z_\ell^\times$. Here the associated extension is finite flat.
One explicit model has components
\[
 \Spec\bigl(\Z_\ell[z]/(z^\ell-u^a)\bigr),\qquad0\leq a<\ell,
\]
with product
\[
 (a,z)(b,w)=
 \left(a+b-\ell\left\lfloor\frac{a+b}{\ell}\right\rfloor,
 \frac{zw}{u^{\lfloor(a+b)/\ell\rfloor}}\right).
\]
The carry identity proves associativity. Its identity is $(0,1)$;
inversion is obtained from the same formula. Every component is finite
free of rank $\ell$, and the resulting group scheme is an extension
of $\Z/\ell$ by $\mu_\ell$ whose fiber above $1$ is the Kummer torsor
of $u$. This gives the required generic-fiber representation.

At two the cyclotomic character is unramified and the inertia image
after twisting is unipotent of order dividing $\ell$. Its restriction
to wild inertia is trivial, since wild inertia is pro-two. The
valuation of the Tate class is $-8$ modulo $\ell$, at an ordinary
point by Lemma~\ref{lem:local-points} and at a cusp by
\eqref{eq:cuspresidual}. It is nonzero. Tame inertia therefore has a
nontrivial transvection, with a one-dimensional invariant subspace.
The Swan conductor is zero and the Artin conductor exponent is one.
\end{proof}

\begin{proposition}[Reducibility]\label{prop:reducible}
For every Selmer section $s$ of $U_\ell$, $\ell\geq5$ prime,
$\overline\rho_s$ is reducible over $\F_\ell$.
\end{proposition}
\begin{proof}
Only finitely many local inertia characters in
Lemma~\ref{lem:local-residual} need be nontrivial. Indeed
$\overline\rho_s$ has finite image and is ramified at finitely many
primes. Away from $2,\ell$, a quadratic scalar cannot be canceled by
a nontrivial unipotent matrix in odd characteristic, so the needed
inertial twist is trivial wherever the original representation is
unramified.

The local quadratic inertia characters patch to a global character
$\eps:G_\Q\to\{\pm1\}$. To see this precisely, use
Kronecker--Weber to identify $G_\Q^{\ab}$ with
$\Zhat^\times=\prod_p\Z_p^\times$. At the prime $p$, the restriction
of a local abelian character to inertia is its character on
$\Z_p^\times$ under local reciprocity. Put the prescribed character
on each of the finitely many relevant factors and the trivial one on
the others. Its restriction at $p$ differs from $\eps_p$ by an
unramified character. Such a difference preserves unramifiedness,
the conductor, and finite flatness; in the last case it is a twist by
an \et ale rank-one group scheme over $\Z_\ell$.

Set $r=\overline\rho_s\otimes\eps$. It has determinant
$\overline\chi_\ell$, is unramified away from $2,\ell$, has
prime-to-$\ell$ conductor $N(r)=2$, and is finite flat at $\ell$.
Serre's finite-flat criterion
\citep[\S2.8, Proposition~4, pp.~189--190]{serre} therefore gives $k(r)=2$.

Suppose $r$ were irreducible over $\F_\ell$. It is absolutely
irreducible. In fact complex conjugation has distinct eigenvalues
$1,-1$ since the determinant is cyclotomic. Any invariant line over
$\overline{\F}_\ell$ must be one of those eigenspaces, both of which
are defined over $\F_\ell$; such a line would contradict
irreducibility over $\F_\ell$.

The refined Serre theorem now produces a weight-two cusp form of
level $N(r)=2$. The modularity input is
\citet[Theorem~9.1, author version p.~18]{kwone}, with the lifting
results proved in \citet[\S10.2--10.3, author version pp.~92--94]{kwtwo}
and hypothesis (H) established by
\citet[Theorem~0.1 and Corollary~0.2, author version p.~2]{kisin}.
For the refinement, Ribet's level-lowering theorem
\citep[Theorem~1.1, p.~432]{ribet} begins with a weight-two modular
representation, whereas Edixhoven's minimum-weight theorem
\citep[Theorem~4.5, author version p.~11]{edixhoven} supplies the
weight component. The combined refinement is Diamond's
Theorem~6.4 \citep{diamond}. We use its precise sufficient formulation
in \citet[Theorem~2.14, pp.~503--504]{wiles}: an odd absolutely
irreducible modular representation in characteristic $\ell>3$,
finite flat at $\ell$ with cyclotomic determinant there, occurs in
weight two at its prime-to-$\ell$ Artin conductor. The theorem
gives a character of order prime to $\ell$; determinant comparison
makes its reduction trivial here, hence the character itself trivial.
Thus the form lies in $S_2(\Gamma_0(2))$. Qualitative modularity
alone would not give this conclusion.

But $\Gamma_1(2)=\Gamma_0(2)$, and the modular curve $X_0(2)$ has
genus zero. Weight-two cusp forms are its holomorphic differentials,
so $S_2(\Gamma_1(2))=0$. This contradiction proves reducibility.
Quadratic scalar twisting does not change invariant lines, so the
same holds for $\overline\rho_s$.
\end{proof}

\section{A modular cover of level \texorpdfstring{$4\ell$}{four times the exponent}}\label{sec:modularcover}

The invariant line now converts representation-theoretic information into
a finite section neighborhood with a useful proper map.

\begin{proposition}\label{prop:fourlevel}
For a Selmer section of $U_\ell$, $\ell\geq5$ prime, there is a
geometrically connected finite \et ale section neighborhood $V\to U_\ell$
whose smooth proper compactification admits a nonconstant morphism
to $X_0(4\ell)$.
\end{proposition}
\begin{proof}
Let $Z_\ell\to T$ be the finite \et ale scheme of lines in the
Legendre $\ell$-torsion. Its geometric fiber is
$\mathbf P^1(\F_\ell)$. An invariant line for
$\overline\rho_s$ is a fixed point of $s(G_\Q)$ in the fiber of
$U_\ell\times_T Z_\ell\to U_\ell$. Its stabilizer is an open
subgroup containing the section, hence defines a geometrically
connected component $V$ to which the section lifts. On $V$ there is
a Legendre elliptic curve $E$ with a tautological finite \et ale
cyclic subgroup $C\subset E[\ell]$ of order $\ell$.

The two sections $P=(0,0)$ and $Q=(1,0)$ are independent points of
$E[2]$. Form $\phi:E\to E'=E/\langle P\rangle$, with dual isogeny
$\phi':E'\to E$. The subgroup
\[
 C_4=(\phi')^{-1}(\langle Q\rangle)
\]
has order four and is cyclic. Indeed, for a geometric lift $R$ of
$Q$ under $\phi'$, the identity $\phi\phi'=[2]_{E'}$ gives
$2R=\phi(Q)\neq0$. Thus $R$ has order four and generates the
whole inverse image. The definition makes $C_4$ Galois stable.

Since $\ell$ is odd, $\phi(C)$ has order $\ell$ and intersects
$C_4$ trivially. Their sum is a cyclic subgroup of $E'$ of order
$4\ell$. The pair $(E',C_4+\phi(C))$ gives a morphism
$V\to Y_0(4\ell)$ to the coarse modular curve. A coarse moduli map
suffices; no choice of a generator or fine moduli interpretation
of $Y_0(4\ell)$ is required.

This map is nonconstant. Otherwise $j(E')$ would be constant. An
elliptic curve over an algebraically closed field has only finitely
many order-two subgroup schemes, so only finitely many two-isogenous
$j$-invariants. The nonconstant function $j(E)=j(\lambda)$ could
not then occur. A morphism of smooth curves into a proper curve
extends across each missing point by the valuative criterion, giving
a nonconstant, hence finite, map $X_V\to X_0(4\ell)$.
\end{proof}

\begin{lemma}\label{lem:modulargenus}
For an odd prime $\ell$,
\[
 g(X_0(4\ell))=\frac{\ell-3}{2}.
\]
\end{lemma}
\begin{proof}
The index of $\Gamma_0(4\ell)$ is
$4\ell(1+1/2)(1+1/\ell)=6(\ell+1)$. The cusp formula
$\sum_{d\mid4\ell}\varphi(\gcd(d,4\ell/d))$ gives six: the six
divisors are $1,2,4,\ell,2\ell,4\ell$, and each summand is one.
There are no elliptic points of order two since $4\mid4\ell$, and
none of order three since $x^2+x+1$ has no root modulo two.
The standard modular genus formula gives
$1+6(\ell+1)/12-6/2=(\ell-3)/2$.
\end{proof}

By \citet[Corollary~8.8, pp.~385--386]{stoll},
\begin{equation}\label{eq:stollmodular}
 X_0(4\ell)(\A_\Q)_\bullet^{\fab}=X_0(4\ell)(\Q)
 \qquad(\ell\geq5).
\end{equation}
This is equality with the rational points, not merely their adelic
closure. Stoll's Definition~6.1(6) is the definition of excellence
used in his corollary. His erratum \citep[\S2(6)]{stollerrata}
preserves the curve statements; its separate real-place correction
is already incorporated in Lemma~\ref{lem:finite-support}.
In particular no unproved assertion about connected components of
twists is inserted in this use of the theorem.

Apply Corollary~\ref{cor:descent-map} to $V$ and the map of
Proposition~\ref{prop:fourlevel}. The compactification of $V$ has
genus at least one because it covers $F_\ell$. Its lifted Selmer
section is a point section or a cusp section, and its image on
$U_\ell$ is of the same kind.

\section{Excluding rational points and treating the cubic}\label{sec:points}

This section completes the prime-exponent proof without assuming Fermat's
Last Theorem as a separate input.

\begin{lemma}[The rational torsion contradiction]\label{lem:torsion}
For every prime $\ell\geq5$, $U_\ell(\Q)=\varnothing$.
\end{lemma}
\begin{proof}
Suppose $x\in U_\ell(\Q)$. Its Legendre curve $E/\Q$ has full
rational two-torsion, and Proposition~\ref{prop:reducible} applies
to its point section. Use the global quadratic twist constructed in
that proposition; since the representation comes from an elliptic
curve in this case, it is the $\ell$-torsion of the corresponding
quadratic twist $E^\eps$.

Let $\alpha,\beta:G_\Q\to\F_\ell^\times$ be the diagonal characters
of its reducible representation. At every prime different from
$\ell$ their inertial restrictions are trivial: this follows from
unramifiedness away from two and from unipotence at two. At $\ell$
the semisimplified inertial characters are $1$ and
$\overline\chi_\ell$. In the potentially multiplicative case this
is the Tate description. In the good ordinary case it is the usual
connected--\et ale description of $E[\ell]$. The good supersingular
case cannot occur in a reducible representation over $\F_\ell$:
its inertia characters have niveau two and have no invariant
$\F_\ell$-line. An unramified quadratic twist leaves these statements
unchanged.

Choose the global diagonal character which is trivial on inertia
at $\ell$. It is unramified at all finite primes. The finite
extension cut out by this character is therefore unramified at all
finite primes of $\Q$ and must be trivial by Minkowski's discriminant
theorem. Hence $\{\alpha,\beta\}=\{1,\overline\chi_\ell\}$.

If the invariant line has character $1$, $E^\eps$ has a rational
point of order $\ell$. If its character is
$\overline\chi_\ell$, its subgroup is $\mu_\ell$; quotienting by it
gives an elliptic curve whose dual isogeny has kernel the Cartier
dual $\Z/\ell$, and thus again a rational point of order $\ell$.
Quadratic twisting acts trivially on two-torsion. An odd-degree
isogeny induces a Galois-equivariant isomorphism on two-torsion.
Consequently the resulting elliptic curve has a rational torsion
subgroup
\[
 (\Z/2)\oplus(\Z/2\ell).
\]
Mazur's theorem permits a subgroup of this form only when
$\ell\leq4$ \citep[Theorem~(8), p.~35]{mazur}.
This contradicts $\ell\geq5$.
\end{proof}

The preceding section and Lemma~\ref{lem:torsion} prove
Theorem~\ref{thm:fermat} for every prime exponent at least five.
For example $X_0(20)$ has genus one and gives exponent five;
$X_0(28)$ has genus two and gives exponent seven. The forbidden
rational torsion groups in these cases are respectively
$(\Z/2)\oplus(\Z/10)$ and $(\Z/2)\oplus(\Z/14)$.

\begin{proposition}[The cubic]\label{prop:cubic}
Every Selmer section of $U_3/\Q$ is cuspidal, and
$F_3(\Q)=\{y_1,y_2,y_3\}$.
\end{proposition}
\begin{proof}
On $X+Y\neq0$ put
\[
 u=-\frac{3Z}{X+Y},\qquad
 v=\frac{9(X-Y)/(X+Y)-1}{2}.
\]
A direct calculation, valid before imposing the Fermat equation, is
\begin{equation}\label{eq:cubicidentity}
 v^2+v-u^3+7=
 \frac{27(X^3+Y^3+Z^3)}{(X+Y)^3}.
\end{equation}
The inverse rational map is
$(X:Y:Z)=(-v-5:v-4:3u)$. Thus the smooth proper curves are isomorphic,
and the cubic is
\[
 E:\quad v^2+v=u^3-7.
\]
Under $x=4u$, $y=8v+4$, this becomes $y^2=x^3-432$.
Its identification with $X_0(27)$ is recorded explicitly in
\citet[\S10.1, p.~65]{krizli}. Therefore Stoll's
Corollary~8.8 applies to $E$.

For clarity, that equality also shows $E(\Q)$ is finite, without
requiring a numerical analytic-rank computation. The finite-descent
set on a proper variety is closed and compact
\citep[discussion preceding Definition~5.4, p.~364]{stoll}. Its projection to
$E(\Q_p)$ is therefore compact, and here it is exactly the injective
image of the countable group $E(\Q)$. An infinite compact metrizable
group is uncountable: otherwise the Baire theorem would give an
isolated point, translations would make it discrete, and compactness
would make it finite. Thus this image is finite.

The discriminant of the displayed integral equation is $-3^9$, so
it has good reduction at five and seven. Counting solutions gives
$\#E(\F_5)=6$ and $\#E(\F_7)=9$. One can do the count by hand from
the discriminant $4u^3-27$ of the quadratic in $v$. At five the
numbers of affine points for $u=0,1,2,3,4$ are $0,0,1,2,2$; at seven
they are $2,0,0,2,0,2,2$. Add the point at infinity in each case.
Prime-to-$p$ rational torsion injects under good reduction at $p$;
using both primes shows that $\#E(\Q)$ divides three.

The three points $O,(3,-5),(3,4)$ are distinct and are the images
of $y_3,y_1,y_2$, respectively. They therefore exhaust $E(\Q)$.
Any Selmer section of $U_3$ has compactification centers in
$E(\A_\Q)_\bullet^{\fab}=E(\Q)$. They lie in this finite set of
cusps. Lemma~\ref{lem:finite-support} proves the assertion.
\end{proof}

\section{Every odd exponent and the three packets}\label{sec:allodd}

Finite \et ale maps now transport cuspidality from prime exponent to any
odd exponent. Let $\ell\mid n$ be prime and put $m=n/\ell$. The map
\begin{equation}\label{eq:powercover}
 U_n\longrightarrow U_\ell,\qquad
 (X:Y:Z)\longmapsto(X^m:Y^m:Z^m)
\end{equation}
is finite \et ale. On the chart $Z=1$ it adjoins $m$th roots of
two invertible functions, and $m$ is invertible in characteristic
zero. Its compactification maps boundary to boundary.

Let $s$ be Selmer on $U_n$. The image section on $U_\ell$ is Selmer,
so by the prime cases it lies in a cusp decomposition group.
Identify $\pi_1(U_n)$ with the open subgroup corresponding to
\eqref{eq:powercover}. Choose the cusp prolongation downstairs so
that its decomposition group contains the image of $s$. The cusp
above it in the normalization has decomposition group equal to the
intersection with this open subgroup. That intersection contains
$s(G_\Q)$, hence maps onto $G_\Q$; its residue field is therefore
$\Q$. Thus it is a rational cusp upstairs and $s$ is cuspidal.

When $n$ is odd, a rational cusp with $X=0$ has $(Y/Z)^n=-1$, so
$Y/Z=-1$. The same argument in the other two charts gives exactly
$y_1,y_2,y_3$. Each has a rational parameter and a nonzero rational
tangent. Stix's Proposition~8 identifies its local section set as a
torsor under $H^1(\Q,\Zhat(1))$. Stix's Theorem~14 makes the maps
from these three local section sets injective with pairwise disjoint
images. Its hypothesis, the injectivity of the proper section map
over finite extensions, holds over number fields by Mordell--Weil,
as explained in Remark~15(2) and Appendix~B of that paper.

Finally every global cusp section is Selmer: restriction to a
completion remains a section in the decomposition group of the same
rational cusp. This proves both inclusions in \eqref{eq:packets},
and completes the proof of Theorem~\ref{thm:fermat}.

\section{From proper local sections to affine sections}\label{sec:local}

This section proves the additional implication needed in
Corollary~\ref{cor:all}; the proper local input is not silently
substituted for an affine assertion.

\begin{lemma}[High-genus neighborhoods]\label{lem:highgenus}
Let $U/k$ be a smooth affine hyperbolic curve in characteristic zero
and let $s$ be a full section. The section neighborhoods whose smooth
proper compactifications have genus at least two are cofinal.
\end{lemma}
\begin{proof}
Let $g$ be the genus of the compactification and let $r$ be its
number of geometric punctures. Each geometric cusp inertia group
is an infinite procyclic subgroup of $\Delta=\pi_1(U_{\kbar})$.
For each puncture choose a finite quotient in which its inertia
image has order greater than an integer $M$; this is possible
because the subgroup is infinite. Intersect the finitely many
kernels and then take a characteristic open subgroup inside the
intersection, using finite generation of $\Delta$. Intersect also
with any preassigned characteristic open subgroup defining a
section neighborhood.

The resulting geometric cover has degree $d$, and inertia orders
$e_i>M$ at the punctures. Since characteristic zero gives tame
ramification, Riemann--Hurwitz reads
\[
 2g'-2=d\left(2g-2+\sum_{i=1}^r(1-e_i^{-1})\right).
\]
Hyperbolicity means $2g-2+r>0$. Choose $M$ so that
$r/M<2g-2+r$; then the displayed quantity is positive.
The arithmetic subgroup $N\,s(G_k)$ defines the desired section
neighborhood with this geometric cover.
\end{proof}

\begin{proposition}[Affine consequence of the proper local theorem]
\label{prop:local-affine}
Assume Hypothesis~\ref{hyp:local} over a fixed finite extension
$k/\Q_p$. Every full section of a smooth affine hyperbolic curve
$U/k$ is a point section or a section at a $k$-rational cusp.
\end{proposition}
\begin{proof}
Use the cofinal neighborhood system of Lemma~\ref{lem:highgenus}.
For each $U_N$, compose its lifted section with
$\pi_1(U_N)\to\pi_1(X_N)$. By the proper local theorem it is the
section of a unique point $x_N\in X_N(k)$.

If $N'\subset N$, the finite map $U_{N'}\to U_N$ extends to
$X_{N'}\to X_N$. Functoriality takes the section of $x_{N'}$ to
the section of $x_N$. The injectivity assertion of the proper local
theorem for $X_N$ implies that $x_{N'}$ maps to $x_N$. Thus the
points form a compatible system. Their image on the original
compactification is a point $x\in X(k)$.

The geometric tower and its Galois descent action are those in the
last paragraph of the proof of Lemma~\ref{lem:finite-support}.
The compatible rational points therefore give a prolongation
fixed by $s(G_k)$, so that $s(G_k)\subset D_{\widetilde x}$.
If $x$ is interior this is its point section; otherwise it is a
section at the rational cusp $x$.
\end{proof}

At a real place the real section theorem \citep[unnumbered Theorem, p.~2]{real} gives the corresponding
point-or-cusp assertion for hyperbolic curves. One can also use its
proper version and Lemma~\ref{lem:highgenus}: every compactified
neighborhood has a nonempty compact real-point set. The inverse
system has a point by compactness, since a common refinement
satisfies each finite collection of compatibility conditions. The
same stabilizer argument applies. Notice that this uses compact
real-point sets, not a choice of a representative preserved by every
map between real connected components.

Applying Proposition~\ref{prop:local-affine} at every finite place
and the real theorem at infinity makes every full section of $U_n$
Selmer. Theorem~\ref{thm:fermat} proves Corollary~\ref{cor:all}.

\section{Modular curves and even exponents}\label{sec:modulareven}

The same finite-support mechanism explains the modular-curve assertion,
whereas the even-exponent assertion is already forced by the real place.

\begin{proof}[Proof of the modular-curve part of Theorem~\ref{thm:other}]
Let $X_1(N)$ have positive genus. Stoll's Corollary~8.8 says
\[
 X_1(N)(\A_\Q)_\bullet^{\fab}=X_1(N)(\Q).
\]
For a Selmer section of $Y_1(N)$, its compactification centers
therefore lie over a single rational point. Apply
Lemma~\ref{lem:finite-support}, with that point as finite support,
to conclude that the section is a point section or a cusp section.

Mazur's Theorem~(7) states directly that all rational points of
$X_1(N)$ are cusps when it has positive genus
\citep[Theorem~(7), p.~35]{mazur}. Equivalently, these are the levels
$N=11$ or $N\geq13$; Theorem~(8) on the same page is the torsion
classification used in Lemma~\ref{lem:torsion}.
The section must consequently be cuspidal at a rational cusp.

Under Hypothesis~\ref{hyp:local},
Proposition~\ref{prop:local-affine} and the real theorem make
every full section Selmer and give the last assertion.
\end{proof}

\begin{proof}[Proof of the even-exponent part of Theorem~\ref{thm:other}]
If $n$ is even, a real solution of
$X^n+Y^n+Z^n=0$ has all three coordinates zero, so
$F_n(\R)=\varnothing$. Neither an interior real point nor a real
cusp exists. This already makes the Selmer condition impossible,
also for any quotient if the local condition is defined as the
image of an interior or cusp section.

Suppose a full section existed. Restrict it to $G_\R$ and take a
high-genus neighborhood over $\R$ as in
Lemma~\ref{lem:highgenus}. Its compactification maps to $F_n$,
so it too has no real point. But its induced proper section and
the real section theorem imply a real point, a contradiction.
This proof covers $n=2$, since $U_2$ is hyperbolic even though
its compactification has genus zero.
\end{proof}

\section{The two corrections to Steklov}\label{sec:corrections}

This section specifies which conclusions are and are not affected by the
two defective steps in the printed proof.

In \citet[Proposition~4.6]{steklov}, the extension at a point of
ramification index $e$ is identified with
$k((j^{-1/e}))/k((j^{-1}))$. A ramified parameter has instead the
form $j^{-1}=a t^e h(t)$, and this identification over $k$ requires
the relevant leading constant to be an $e$th power. It does not
follow from the ramification index. Proposition~\ref{prop:cusp-class}
gives $a=2^{-8}$ for the Fermat cusps. Unit tangent gives a class
whose two-adic valuation is $-8$, so it is not divisible by any
odd exponent $n\geq3$.

Accordingly the split residual description used in the proof of
Steklov's Proposition~5.10 is false. Its conclusion that cuspidal
sections have the required local Fermat conditions remains true:
Lemma~\ref{lem:local-residual} proves unramifiedness at odd
$p\neq\ell$, finite flatness at $\ell$, and tame conductor exponent
one at two. These imply the conductor bound in his Definition~5.8.
With that proposition reproved, the conclusion for
$\ell\geq11$ stands. Our proof does not assume that conclusion.

The second error is independent. The last case of his Theorem~5.14
claims that $Y_1(\ell)(\Q)$ is empty for $\ell\geq7$ by Mazur's
torsion theorem. Mazur's list includes $\Z/7$. An explicit example is
\[
 E:\quad y^2-xy-4y=x^3-4x^2,
 \qquad P=(0,0).
\]
Its discriminant is $-1664\neq0$, and the chord-and-tangent law gives
\[
 2P=(4,8),\quad3P=(2,2),\quad4P=(2,4),\quad
 5P=(4,0),\quad6P=(0,4)=-P.
\]
Hence $P$ has exact order seven. The printed proof does not establish
exponent seven. Lemma~\ref{lem:torsion} instead preserves full
two-torsion and excludes $(\Z/2)\oplus(\Z/14)$, which is absent
from Mazur's list. The same argument excludes
$(\Z/2)\oplus(\Z/10)$ for exponent five.

\section{Scope and further questions}\label{sec:scope}

The global full-group Selmer theorem is proved for all odd exponents.
Removing its local hypothesis requires the proper local input in
Hypothesis~\ref{hyp:local}; the reduction from that precise statement
is proved in Section~\ref{sec:local}. The unrefereed status of the
asserted input is not changed by that deduction.

For quotients of the geometric fundamental group, the modular
representation used here need not factor through the quotient.
In particular a full-group proof does not automatically prove a
prosolvable, metabelian, or pro-$3$ theorem. Those questions, their
proved reductions, and the distinction between arbitrary
noncuspidal sections and Selmer sections are treated separately
in the companion paper. No theorem for a smaller group is inferred
merely by projecting the full section set.

\subsection*{Method and records}

The mathematics and text were produced by GPT-6 Astra Pro in ChatGPT
Pro mode. Separate instances of the same model attempted the main
statements from their statements before comparing proofs; this is
same-model checking, not independent human verification. Earlier
global arguments were examined by Claude Opus 5.5 at maximal effort
on 25 November 2026, and selected earlier computations by that model
at high effort on 24 November 2026. Those checks have their stated
scope and do not certify the local preprint or every extension to
smaller groups. No human has checked this paper.

The original text, arguments, and programs are dedicated to the public
domain under \href{https://creativecommons.org/publicdomain/zero/1.0/}{CC0~1.0}.
The program \emph{Fermat arithmetic checks}, its expected output,
and the detailed statement-by-statement checks are in the
\href{\workbenchurl}{workbench}. The permanent address of the
records is \url{\recordurl}. Published sources are cited there and
are not redistributed.

\IfFileExists{identity.tex}{\input{identity.tex}}{}

\begin{thebibliography}{99}
\bibitem[Betts et~al.(2026)]{bkl}
Betts, L.~A., Kumpitsch, T., \& L\"udtke, M. (2026).
\emph{Chabauty--Kim, finite descent, and the Section Conjecture for locally
geometric sections}. arXiv:2305.09462v2.
\url{https://arxiv.org/abs/2305.09462v2}.

\bibitem[Bresciani and Vistoli(2020)]{real}
Bresciani, G., \& Vistoli, A. (2020).
\emph{An elementary approach to Stix's proof of the real section conjecture}.
arXiv:2012.06278v2, unnumbered Theorem, p.~2.
\url{https://arxiv.org/abs/2012.06278v2}.

\bibitem[Diamond(1995)]{diamond}
Diamond, F. (1995). The refined conjecture of Serre.
In J. Coates \& S.-T. Yau (Eds.),
\emph{Elliptic curves, modular forms, \& Fermat's last theorem}
(pp.~22--37). International Press.

\bibitem[Edixhoven(1992)]{edixhoven}
Edixhoven, B. (1992). The weight in Serre's conjectures on modular forms.
\emph{Inventiones Mathematicae, 109}, 563--594.
\url{https://doi.org/10.1007/BF01232041}.

\bibitem[Grothendieck(1981)]{grothendieck}
Grothendieck, A. (1981). \emph{La Longue Marche \`a travers la th\'eorie
de Galois}, \S51.16. Manuscript, Archives Grothendieck, Universit\'e de
Montpellier, cote~140-4, fols.~733--736v.
\url{https://grothendieck.umontpellier.fr/140-4.pdf}.

\bibitem[Khare and Wintenberger(2009a)]{kwone}
Khare, C., \& Wintenberger, J.-P. (2009a).
Serre's modularity conjecture (I).
\emph{Inventiones Mathematicae, 178}, 485--504.
\url{https://doi.org/10.1007/s00222-009-0205-7}.

\bibitem[Khare and Wintenberger(2009b)]{kwtwo}
Khare, C., \& Wintenberger, J.-P. (2009b).
Serre's modularity conjecture (II).
\emph{Inventiones Mathematicae, 178}, 505--586.
\url{https://doi.org/10.1007/s00222-009-0206-6}.

\bibitem[Kisin(2009)]{kisin}
Kisin, M. (2009). Modularity of 2-adic Barsotti--Tate representations.
\emph{Inventiones Mathematicae, 178}, 587--634.
\url{https://doi.org/10.1007/s00222-009-0207-5}.

\bibitem[Kriz and Li(2019)]{krizli}
Kriz, D., \& Li, C. (2019). Goldfeld's conjecture and congruences between
Heegner points. \emph{Forum of Mathematics, Sigma, 7}, e15.
\url{https://doi.org/10.1017/fms.2019.9}.

\bibitem[Mazur(1977)]{mazur}
Mazur, B. (1977). Modular curves and the Eisenstein ideal.
\emph{Publications Math\'ematiques de l'IH\'ES, 47}, 33--186.
\url{https://doi.org/10.1007/BF02684339}.

\bibitem[OpenAI(2026)]{openai}
OpenAI. (2026, October 6). \emph{The $p$-adic section conjecture}.
Unrefereed preprint, Theorem~1.1, p.~2; first version September 24.
\url{https://github.com/openai/math}.

\bibitem[Porowski(2025)]{porowski}
Porowski, W. (2025). \emph{On reductions of Selmer sections}.
arXiv:2509.15660v1.
\url{https://arxiv.org/abs/2509.15660v1}.

\bibitem[Ribet(1990)]{ribet}
Ribet, K.~A. (1990). On modular representations of
$\Gal(\overline{\Q}/\Q)$ arising from modular forms.
\emph{Inventiones Mathematicae, 100}, 431--476.
\url{https://doi.org/10.1007/BF01231195}.

\bibitem[Serre(1987)]{serre}
Serre, J.-P. (1987). Sur les repr\'esentations modulaires de degr\'e~2
de $\Gal(\overline{\Q}/\Q)$.
\emph{Duke Mathematical Journal, 54}, 179--230.
\url{https://doi.org/10.1215/S0012-7094-87-05413-5}.

\bibitem[Steklov(2026)]{steklov}
Steklov, B. (2026, March 3). \emph{Fermat's Last Theorem for Selmer
sections}. arXiv:2602.17746v2.
\url{https://arxiv.org/abs/2602.17746v2}.

\bibitem[Stix(2008)]{stix}
Stix, J. (2008). \emph{On cuspidal sections of algebraic fundamental
groups}. arXiv:0809.0017; author PDF dated November 13, 2018.
\url{https://arxiv.org/abs/0809.0017}.

\bibitem[Stoll(2007)]{stoll}
Stoll, M. (2007). Finite descent obstructions and rational points on
curves. \emph{Algebra \& Number Theory, 1}(4), 349--391.
\url{https://doi.org/10.2140/ant.2007.1.349}.

\bibitem[Stoll(2017)]{stollerrata}
Stoll, M. (2017, March 29). \emph{Errata: Finite descent obstructions
and rational points on curves}.
\url{https://msp.org/ant/2007/1-4/ant-v1-n4-x01-Errata-FiniteDescent.pdf}.

\bibitem[Wiles(1995)]{wiles}
Wiles, A. (1995). Modular elliptic curves and Fermat's Last Theorem.
\emph{Annals of Mathematics, 141}(3), 443--551.
\url{https://doi.org/10.2307/2118559}.
\end{thebibliography}
\end{document}
